% Source: https://artofproblemsolving.com/community/c3h2191161_this_forum_is_for_number_theory_doubt_solving
% Stable identity: aops-topic-2191161-q01
\textbf{Existence.} By the division algorithm there are integers $q_0$ and $r_0$ with $a=bq_0+r_0$ and $0\le r_0<b$. Put $q=q_0-2$ and $r=r_0+2b$. Then
\[bq+r=bq_0-2b+r_0+2b=a\qquad\text{and}\qquad 2b\le r_0+2b<3b .\]

\textbf{Uniqueness.} Suppose $a=bq+r=bq'+r'$ with $2b\le r<3b$ and $2b\le r'<3b$. Then $b(q-q')=r'-r$. Since both $r$ and $r'$ lie in the interval $[2b,3b)$ of length $b$, we have $|r'-r|<b$. So $|q-q'|\cdot b<b$, which forces $|q-q'|<1$, that is, $q=q'$; then $r=r'$ as well. $\blacksquare$

(In fact $q=\left\lfloor\dfrac {a}{b}\right\rfloor-2$ and $r=a-bq$.)
